% \noindent
% The goal of this thesis was to use machine-learning techniques to automatically classify security-sensitive Java methods using Javadoc. \\

% \noindent
% A dataset of security-relevant Java methods' source-code and Javadoc was produced. To our knowledge, a similar dataset does not already exist in the literature. The final dataset (training and validation combined) is comprised of 1096 methods (158 sources, 268 sinks, 175 sanitisers, and 495 security-insensitive methods). The dataset is produced dynamically from a preliminary dataset format. The preliminary dataset and the code to dynamically build the dataset is publicly available on \href{https://github.com/uqcyber/SecurityMethodClassifier/tree/main/tools/dataset-manager}{GitHub}. \\

% \noindent
% In addition to the training and validation datasets, a test dataset was created on a more realistic use-case. This test dataset contained a random sample of methods from Apache's Log4j ($n = 165$) which were manually classified. This test dataset was used to evaluate the models' performance.  \\

% \noindent
% The pre-trained models CodeBERT and UniXcoder were fine-tuned on our dataset and used to perform the classification. Both mean and CLS pooling strategies were used as it was not clear which was more effective for this use-case. It was found that CLS pooling performed better for UniXcoder and mean pooling performed better for CodeBERT. \\

% \noindent
% The highest performing model configuration was found to be fine-tuned CodeBERT with a neural network classifier using mean pooling. This model achieved an F1-score of $0.68$. This result is significantly better than the performance of a naive SVM classifier using tf-idf for feature generation, which had an F1-score of 0.51. \\

% % \noindent
% % To our knowledge, the approach we present in this thesis is novel. Furthermore, the results of this thesis indicate that Javadoc contains information useful for security analyses.

% \noindent
% The approach we have taken in this thesis, to our knowledge, is novel. While the models we have presented in this thesis are not accurate enough to allow for taint analysis configurations to be fully automatically generated, as was hoped at beginning of this thesis, the models do perform well enough to allow for security analysts to create configurations more efficiently than from scratch. Furthermore, it is possible our approach could be combined with existing approaches in the literature to produce an even more performant model. Most importantly, however, this thesis has shown that Javadoc does indeed contain information useful for security anaylsis and it is worth exploring what other security analyses or other tasks could be enhanced using deep learning techniques on Javadoc.

\noindent
In this concluding chapter of the thesis, we summarise the project's outcomes and discuss potential future work for this project.

\section{Project Outcomes}

\noindent
A dataset of security-relevant Java methods' source code and Javadoc was produced. To our knowledge, a similar dataset does not already exist in the literature. The final dataset (training and validation combined) is comprised of 1096 methods (158 sources, 268 sinks, 175 sanitisers, and 495 security-insensitive methods). The dataset is produced dynamically from a preliminary dataset. The preliminary dataset and the code to dynamically build the dataset is publicly available on \href{https://github.com/uqcyber/SecurityMethodClassifier/tree/main/tools/dataset-manager}{GitHub}. \\

\noindent
In addition to this dataset, a test dataset was created. This test dataset consisted of a random sample of Java methods from Apache's Log4J that were manually labelled. This test dataset served as the benchmark test for our models as it more accurately reflects a realistic use case for the models. \\

\noindent
Furthermore, several machine learning models based on CodeBERT and UniXcoder were produced. These pre-trained models were evaluated and tested in a variety of configurations. The difference between different pooling strategies --- mean versus CLS pooling --- was investigated. It was determined that UniXcoder performed best with CLS pooling, and CodeBERT performed best with mean pooling. \\

\noindent
The highest-performing model produced for this thesis project was CodeBERT with a neural-network classifier using mean pooling. This model achieved an F1-score of $0.68$ on the test dataset. The source code for all the machine-learning models and configurations produced for this thesis can be found on \href{https://github.com/uqcyber/SecurityMethodClassifier/tree/main/models}{GitHub}. \\


\section{Future Work}

\noindent
In addition to the improvements suggested in Section \ref{limitations_and_improvements}, there are a number of extensions to this project that can be undertaken. \\

\noindent
One possible extension is to include the \emph{context} of a method call into the training data. For instance, it may be useful to include the Javadoc for the method that calls the method to be classified into the training data. This is likely to be a difficult extension to implement as it would require a call-graph analysis to be conducted. \\ 

\noindent
Another potential extension would be to include in-line comments in the training data. So far, in-line comments have been removed, but perhaps further performance could be gained by including all the in-line comments of a method (in addition to the Javadoc) in the training data. Recall that the dataset provided for this thesis is not static; it is produced dynamically by the dataset manager. As such, implementing this extension would likely only require small tweaks to the dataset manager and for the models to be trained once again. \\

\noindent
Another simple extension is to `highlight' keywords in the input during training and ensure the tokeniser does not split them. Furthermore, these tokens could have their embeddings fixed so their embeddings are not altered by contextual information. It is not clear at this stage what the keywords would be, but it is worth exploring this idea in future.\\

\noindent
Finally, learning could be conducted on an additional classification task where the model learns to classify each token in the input as either security-sensitive or security-insensitive. This approach was inspired by Wang et. al's identifier-tagging approach \cite{identifier_tagging} and may allow for more sophisticated learning to occur.